\documentclass{article}

% -- PACKAGES --
\usepackage[utf8]{inputenc}
\usepackage[
    backend=biber,
    style=authoryear,
    maxcitenames=2, % "Author & Author" for 2, "Author et al." for 3+ in text
    maxbibnames=8,  % Truncate to et al. in the bibliography if more than 8 authors
    uniquelist=false,
    uniquename=false,
    giveninits=true % Uses initials for first names
]{biblatex}
\usepackage{amsmath}
\usepackage{amssymb}
\usepackage{mathtools}
\usepackage{hyperref}
\usepackage{graphicx}
\usepackage{float}
\usepackage{subcaption}
\usepackage{microtype}
\usepackage[dvipsnames]{xcolor}
\usepackage[skip=5pt]{parskip}
\usepackage[left=2.5cm, right=2.5cm, top=2.5cm, bottom=2.5cm]{geometry}
% -------------

% -- PACKAGE CONFIGURATIONS --
\addbibresource{references.bib} % Link to .bib file

% Variables
\newcommand{\plotsize}{0.6}
% ----------------------------

% Title features
\title{ASP3231 Assignment 3}
\author{Alec Stanley\\
\vspace{0.1cm}
\small ID: 33963975}
\date{\today}

\begin{document}

\maketitle

\section*{Question 1}

\subsection*{1a.}
Wien's Displacement Law relates the wavelength at which a blackbody is most radiant to the inverse of its temperature. It is defined under the assumption that the object is a perfectly radiating blackbody, which is a typical assumption for analysing most stars. It follows that

\begin{equation}
    \lambda_\textup{max}=\frac{2.898\cdot10^{-3}\textup{m}\cdot\textup{K}}T.
\end{equation}

On the t=1.43 days graph in Figure 1, the wavelength that is emitted at the highest flux (intensity) is approximately $1.2\cdot10^{3}$ angstroms, or $1.2\cdot10^{-7}$ metres. Hence, the temperature of the object 1.43 days after the explosion is $\approx2.415\cdot10^{4}$ K.

\subsection*{1b.}

\subsection*{1c.}

\subsection*{1d.}

\subsection*{1e.}

\section*{Question 2}

\subsection*{2a.}

\subsection*{2b.}

\subsection*{2c.}

\subsection*{2d.}

\printbibliography[heading=bibnumbered]

\end{document}
